\slajdid{03-s011}
\begin{frame}[label=03-s011]
\gusttekst
Takođe ulazi i izlazi mogu biti i sa aktivnim nivoima logičke nule. Na primer da su izlazi sa aktivnim logičkim nulama
\begin{columns}[c,onlytextwidth]
\begin{column}{.55\textwidth}\centering
\begingroup
% Ujednačen mali simbol sa 8 pt oznakama; selekcioni priključci su levo.
% \DEdek{ime}{x}{y}; globalno DEbar/DEbubble menjaju izlaze.
\providecommand{\DEbar}{0}
\providecommand{\DEbubble}{0}
\newcommand{\DEdek}[3]{%
\begin{scope}[shift={(#2,#3)},font=\fontsize{8}{8.5}\selectfont]
\draw[DEvod] (0,0) rectangle (1.7,1.65);
\node[align=center,inner sep=1pt] at (.58,1.27) {Dekoder\\2/4};
\foreach \n/\y in {3/1.24,2/.96,1/.68,0/.40}{
 \ifnum\DEbar=1
  \node[anchor=east,inner sep=1pt] at (1.7,\y) {$\overline{\mathrm{Y\n}}$};
 \else
  \node[anchor=east,inner sep=1pt] at (1.7,\y) {Y\n};
 \fi
 \ifnum\DEbubble=1
  \draw[DEvod,fill=white] (1.75,\y) circle(.05);
  \draw[DEvod] (1.80,\y)--(1.95,\y);
 \else
  \draw[DEvod] (1.7,\y)--(1.95,\y);
 \fi
 \coordinate (#1-y\n) at (1.95,\y);
}
\foreach \lab/\yy in {S1/.96,S0/.68,CS/.40}{
 \node[anchor=west,inner sep=1pt] at (0,\yy) {\lab};
 \draw[DEvod] (-.2,\yy)--(0,\yy);
 \coordinate (#1-\lab) at (-.2,\yy);
}
\end{scope}}

\begin{tikzpicture}[x=1cm,y=1cm,font=\fontsize{9}{10}\selectfont]
\def\DEbar{1}\def\DEbubble{0}\DEdek{a}{0}{0}
\node at (2.3,.83) {$=$};
\def\DEbar{0}\def\DEbubble{1}\DEdek{b}{2.9}{0}
\node at (5.2,.83) {$\ne$};
\def\DEbar{1}\def\DEbubble{1}\DEdek{c}{5.8}{0}
\end{tikzpicture}
\endgroup

\end{column}
\begin{column}{.42\textwidth}
Obratite pažnju da su i prvi i drugi način označavanja isti, i mogu slobodno da se koriste, dok je treći pogrešan i takva komponenta nije isto što i prve dve. Suštinski treća komponenta bi opet bila dekoder sa aktivnim logičkim jedinicama na izlazu, postoji dvostruka negacija. Način prikazan na drugom simbolu se i najčešće koristi.
\end{column}
\end{columns}
Primer sa više selekcionih signala.
\par\centering
\begin{tikzpicture}[x=1cm,y=1cm,font=\fontsize{8}{9}\selectfont]
\foreach \i/\x in {0/0,1/3,2/6}{
 \begin{scope}[shift={(\x,0)}]
 \draw[DEvod] (0,0) rectangle(1.65,2.25);
 \node[align=center] at (.82,1.9) {Dekoder\\2/4};
 \foreach \n/\y in {3/1.48,2/1.16,1/.84,0/.52}{
  \node[anchor=east,inner sep=2pt] at (1.65,\y) {Y\n};
  \draw[DEvod] (1.65,\y)--(1.85,\y);
 }
 \foreach \n/\xx in {1/.48,0/1.18}{
  \node[anchor=south,inner sep=2pt] at (\xx,0) {S\n};
  \draw[DEvod] (\xx,0)--(\xx,-.2);
 }
 \node[anchor=west,inner sep=2pt] at (0,.47) {CS1};
 \draw[DEvod] (-.2,.47)--(0,.47);
 \ifnum\i=1
  \node[anchor=west,inner sep=2pt] at (0,.82) {CS2};
 \else
  \node[anchor=west,inner sep=2pt] at (0,.82) {$\overline{\mathrm{CS2}}$};
 \fi
 \ifnum\i=0
  \draw[DEvod] (-.2,.82)--(0,.82);
 \else
  \draw[DEvod,fill=white] (-.05,.82) circle(.05);
  \draw[DEvod] (-.2,.82)--(-.1,.82);
 \fi
 \end{scope}
}
\node[font=\fontsize{11}{12}\selectfont] at (2.4,1.1) {$=$};
\node[font=\fontsize{11}{12}\selectfont] at (5.4,1.1) {$\ne$};
\end{tikzpicture}
\par
Unutrašnji selekcioni signal se formira kao \(CS=CS1\cdot\overline{CS2}\).
\end{frame}
